%! AP Statistics — Unit 1, Lesson 1.9 — Exit Ticket ANSWER KEY
\documentclass[10pt]{article}
\usepackage{apstats-article}
\usepackage{apstats-key}

%=============================================================================
\begin{document}

\pageheader{Unit 1, Lesson 1.9}{Exit Ticket --- Answer Key}
\begin{tabularx}{\linewidth}{X X X}
Name:\;\ans{Key} & Date:\; & Period:\;
\end{tabularx}

\vspace{0.18in}

\begin{tcolorbox}[colback=sky, colframe=navy, boxrule=0.9pt, arc=2mm,
    left=4mm, right=4mm, top=3mm, bottom=3mm]
\small A meteorologist recorded monthly rainfall totals (inches) for two cities
over a year. The parallel boxplots below summarize the distributions.
City B has one high outlier.

\vspace{0.14in}
\begin{center}
\begin{tikzpicture}[scale=1.0]
  \draw[thick] (0.0, 0) -- (7.50, 0);
  \foreach \v/\x in {0/0.00, 1/0.60, 2/1.20, 3/1.80, 4/2.40, 5/3.00,
                      6/3.60, 7/4.20, 8/4.80, 9/5.40, 10/6.00, 11/6.60, 12/7.20} {
    \draw (\x, -0.12) -- (\x, 0.12);
    \node[below, font=\small] at (\x, -0.12) {\v};
  }
  \node[below, font=\small\itshape] at (3.60, -0.56) {Monthly rainfall (inches)};
  \node[left, font=\small] at (0.0, 1.10) {City A};
  \draw[thick, navy] (1.20, 1.10) -- (2.40, 1.10);
  \draw[thick, navy] (2.40, 0.85) rectangle (4.80, 1.35);
  \draw[thick, navy] (3.60, 0.85) -- (3.60, 1.35);
  \draw[thick, navy] (4.80, 1.10) -- (6.60, 1.10);
  \draw[thick, navy] (1.20, 0.95) -- (1.20, 1.25);
  \draw[thick, navy] (6.60, 0.95) -- (6.60, 1.25);
  \node[left, font=\small] at (0.0, 0.35) {City B};
  \draw[thick, navy] (0.00, 0.35) -- (0.60, 0.35);
  \draw[thick, navy] (0.60, 0.10) rectangle (2.40, 0.60);
  \draw[thick, navy] (1.20, 0.10) -- (1.20, 0.60);
  \draw[thick, navy] (2.40, 0.35) -- (4.80, 0.35);
  \draw[thick, navy] (0.00, 0.20) -- (0.00, 0.50);
  \draw[thick, navy] (4.80, 0.20) -- (4.80, 0.50);
  \filldraw[navy] (6.00, 0.35) circle (0.11);
\end{tikzpicture}
\end{center}
\end{tcolorbox}

\vspace{0.16in}

\begin{enumerate}[leftmargin=1.5em, itemsep=8pt]

  \item \renewcommand{\arraystretch}{1.5}
        \begin{tabularx}{\linewidth}{@{} l Y Y Y Y Y @{}}
        \rowcolor{sky}
        & $\min$ & $Q_1$ & $M$ & $Q_3$ & right whisker end \\
        \midrule
        City A & \ans{2} & \ans{4} & \ans{6} & \ans{8} & \ans{11} \\
        City B & \ans{0} & \ans{1} & \ans{2} & \ans{4} & \ans{8}  \\
        \end{tabularx}\\[5pt]
        City A IQR $= 8 - 4 = $ \ans{4 inches} \hfill
        City B IQR $= 4 - 1 = $ \ans{3 inches}

  \item Upper fence $= 4 + 1.5(3) = 4 + 4.5 = $ \ans{8.5 inches}\\[4pt]
        Is 10 an outlier? \ans{Yes} \quad Because: \ans{$10 > 8.5$ (upper fence)}

  \item \ans{City A has a higher median (6 inches) than City B (2 inches), by
        4 inches per month.}

  \item \ans{City B's IQR (3 inches) is \textit{smaller than} City A's IQR
        (4 inches), meaning City B's monthly rainfall is slightly more
        consistent in the middle 50\% of months. However, City B has a high
        outlier at 10 inches, indicating at least one unusually wet month
        compared to its typical range.}

\end{enumerate}

\vspace{0.16in}

\begin{tcolorbox}[colback=redbg, colframe=redacc, boxrule=0.8pt, arc=2mm,
    left=4mm, right=4mm, top=3mm, bottom=3mm]
\small\textbf{AP-Style Practice --- Model Response}\\[6pt]
\ans{The distribution of monthly rainfall for City A is approximately symmetric,
while City B's distribution is right-skewed due to a high outlier at 10 inches.
City B's outlier exceeds the upper fence of 8.5 inches ($= 4 + 1.5 \times 3$);
City A has no outliers. The median monthly rainfall for City A (6 inches) was
\textit{4 inches higher than} the median for City B (2 inches), indicating that
City A is generally wetter. The IQR of monthly rainfall for City A (4 inches) was
\textit{slightly larger than} City B's IQR (3 inches), so City A's rainfall is
somewhat less consistent month to month among its middle 50\% of months.}
\end{tcolorbox}

\begin{teachernote}
\textbf{Grading notes (8 pts total):}
\begin{itemize}[itemsep=2pt]
  \item Q1: 2 pts --- all 10 five-number values correct (1 pt each city) + both IQRs (part of Q1 row).
  \item Q2: 2 pts --- fence shown (1 pt) + correct outlier conclusion (1 pt).
  \item Q3: 1 pt --- states City A higher, names both cities, gives difference (4 in).
  \item Q4: 1 pt --- compares IQRs with direction word; references context (monthly rainfall).
  \item AP-Style: 2 pts --- 1 pt for center comparison (both cities, comparative word, values);
        1 pt for spread comparison (IQRs, both named, context). Shape/outlier addressed = bonus.
\end{itemize}
\end{teachernote}

\end{document}
